%-------------------------
% Resume in Latex
% Author : Abey George
% Based off of: https://github.com/sb2nov/resume
% License : MIT
%------------------------


\documentclass[letterpaper,11pt]{article}


\usepackage{latexsym}
\usepackage[empty]{fullpage}
\usepackage{titlesec}
\usepackage{marvosym}
\usepackage[usenames,dvipsnames]{color}
\usepackage{verbatim}
\usepackage{enumitem}
\usepackage[hidelinks]{hyperref}
\usepackage[english]{babel}
\usepackage{tabularx}
\usepackage{fontawesome5}
\usepackage{multicol}
\usepackage{graphicx}
\usepackage{needspace}
\setlength{\multicolsep}{-3.0pt}
\setlength{\columnsep}{-1pt}
\input{glyphtounicode}


\RequirePackage{tikz}
\RequirePackage{xcolor}
\RequirePackage{fontawesome}
\usepackage{tikz}
\usetikzlibrary{svg.path}




\definecolor{cvblue}{HTML}{0E5484}
\definecolor{black}{HTML}{130810}
\definecolor{darkcolor}{HTML}{0F4539}
\definecolor{cvgreen}{HTML}{3BD80D}
\definecolor{taggreen}{HTML}{00E278}
\definecolor{SlateGrey}{HTML}{2E2E2E}
\definecolor{LightGrey}{HTML}{666666}
\colorlet{name}{cvblue}
\colorlet{tagline}{cvblue}
\colorlet{heading}{cvblue}
\colorlet{headingrule}{cvblue}
\colorlet{accent}{darkcolor}
\colorlet{emphasis}{SlateGrey}
\colorlet{body}{LightGrey}






%----------FONT OPTIONS----------
% Calibri-metric-compatible (free, ships with TeX Live).
\usepackage[sfdefault]{carlito}

% Fallback options (uncomment one if Carlito is unavailable):
% \usepackage{helvet}\renewcommand{\familydefault}{\sfdefault}
% \usepackage[default]{sourcesanspro}
% \usepackage[sfdefault]{FiraSans}


% Adjust margins
\addtolength{\oddsidemargin}{-0.6in}
\addtolength{\evensidemargin}{-0.5in}
\addtolength{\textwidth}{1.19in}
\addtolength{\topmargin}{-.7in}
\addtolength{\textheight}{1.4in}


\urlstyle{same}


\raggedbottom
\raggedright
\setlength{\tabcolsep}{0in}


% Sections formatting (clean corporate look: bold, Title Case, thin rule)
\titleformat{\section}{
  \vspace{-4pt}\raggedright\large\color{heading}\bfseries
}{}{0em}{}[\color{headingrule}\titlerule \vspace{-5pt}]


% Ensure that generate pdf is machine readable/ATS parsable
\pdfgentounicode=1


%-------------------------
% Custom commands
\newcommand{\resumeItem}[1]{
  \item\small{
    {#1 \vspace{-2pt}}
  }
}


\newcommand{\classesList}[4]{
    \item\small{
        {#1 #2 #3 #4 \vspace{-2pt}}
  }
}


\newcommand{\resumeSubheading}[4]{
  \vspace{-2pt}\item
    \begin{tabular*}{1.0\textwidth}[t]{l@{\extracolsep{\fill}}r}
      \textbf{\large#1} & \textbf{\small #2} \\
      \textit{\large#3} & \textit{\small #4} \\
     
    \end{tabular*}\vspace{-7pt}
}


\newcommand{\resumeSubSubheading}[2]{
    \item
    \begin{tabular*}{0.97\textwidth}{l@{\extracolsep{\fill}}r}
      \textit{\small#1} & \textit{\small #2} \\
    \end{tabular*}\vspace{-7pt}
}




\newcommand{\resumeProjectHeading}[2]{
    \item
    \begin{tabular*}{1.001\textwidth}{l@{\extracolsep{\fill}}r}
      \small#1 & \textbf{\small #2}\\
    \end{tabular*}\vspace{-7pt}
}


\newcommand{\resumeSubItem}[1]{\resumeItem{#1}\vspace{-4pt}}


\renewcommand\labelitemi{$\vcenter{\hbox{\tiny$\bullet$}}$}
\renewcommand\labelitemii{$\vcenter{\hbox{\tiny$\bullet$}}$}


\newcommand{\resumeSubHeadingListStart}{\begin{itemize}[leftmargin=0.0in, label={}]}
\newcommand{\resumeSubHeadingListEnd}{\end{itemize}}
\newcommand{\resumeItemListStart}{\begin{itemize}[topsep=2pt, partopsep=0pt, parsep=0pt, itemsep=1pt, leftmargin=0.15in]}
\newcommand{\resumeItemListEnd}{\end{itemize}\vspace{-5pt}}




\newcommand\sbullet[1][.5]{\mathbin{\vcenter{\hbox{\scalebox{#1}{$\bullet$}}}}}


%-------------------------------------------
%%%%%%  RESUME STARTS HERE  %%%%%%%%%%%%%%%%%%%%%%%%%%%%




\begin{document}


%----------HEADING----------




\begin{center}
    {\Huge \bfseries \color{name} Somasekhar Agoli} \\[3pt]
    {\large \color{tagline} Network Automation Engineer} \\[5pt]
    \small
    \raisebox{-0.1\height}\faMapMarker\ Bangalore, KA, India ~
    \href{tel:+916302364851}{\raisebox{-0.1\height}\faPhone\ \underline{+91-6302364851}} ~
    \href{mailto:somasekharagolie1@gmail.com}{\raisebox{-0.2\height}\faEnvelope\ \underline{somasekharagolie1@gmail.com}} ~
    \href{https://www.linkedin.com/in/somasekhar-agoli/}{\raisebox{-0.2\height}\faLinkedinSquare\ \underline{linkedin.com/in/somasekhar-agoli}}
\end{center}


\section{Summary}
\textbf{Network Automation Engineer} with \textbf{2.5+ years at Cisco} engineering \textbf{Python / pyATS} automation for production \textbf{Cisco Secure Firewall infrastructure} (FMC, FTD, cdFMC). Specialized in \textbf{HA failover, disaster recovery, backup/restore, and cross-platform migration} of network security platforms, with deep \textbf{data-plane validation} via traffic generation, event correlation, and backend rule inspection across inline, switched, and routed interface modes (IPv4/IPv6). Authored \textbf{500+ automated test cases} integrated into release \textbf{CI/CD} pipelines that have cut manual validation and recovery effort by up to \textbf{90\%}, pairing hands-on Cisco firewall validation with strong \textbf{networking foundations} (TCP/IP, Routing \& Switching, NAT, ACLs, SSL/TLS).

%-----------CORE COMPETENCIES-----------
\section{Core Competencies}
\small
Cisco Secure Firewall Automation \& Validation (FMC / FTD / cdFMC)~$\bullet$~Networking Foundations~$\bullet$~HA Failover \& Service-Continuity Validation~$\bullet$~Data-Plane Traffic Generation \& Event Validation (IPv4/IPv6)~$\bullet$~Disaster Recovery \& Backup/Restore Automation~$\bullet$~Firewall Migration \& Multi-Mode Interface Validation (Inline / Switched / Routed)~$\bullet$~Python / pyATS Test Framework Engineering~$\bullet$~REST API \& CLI Device Orchestration~$\bullet$~CI/CD Release-Gate Integration (Jenkins)
\vspace{4pt}




%-----------Experience-----------
\section{Experience}

\resumeSubHeadingListStart

\resumeSubheading
  {Cisco Systems}{Nov 2023 -- Present}
  {Network Automation Engineer}{Bangalore, KA, India} 

\needspace{8\baselineskip}
\textbf{Project 1: FMC Migration \& Data-Plane Stress Validation -- Multi-Mode Firewall Topologies}
\nopagebreak
\resumeItemListStart
  \resumeItem{Designed and built a \textbf{pyATS migration validation framework} automating the full lifecycle -- source preparation, bundle upload, migration execution, and post-migration validation of configuration persistence, licensing state, and HA re-formation -- across \textbf{standalone, HA, and dual-instance} topologies spanning \textbf{VMware, AWS, GCP, and physical firewall appliances}.}
  \resumeItem{Built a \textbf{pre-migration heavy-load test harness} -- \textbf{49 test cases} that provision heavy \textbf{Access Control, Prefilter, NAT, SSL/TLS, Security Intelligence, and Network Discovery} policies plus \textbf{Inline Sets, Bridge Groups, Routed Interfaces, Static Routes, and Security Zones} -- to stress the FTD before migration.}
  \resumeItem{Automated \textbf{end-to-end data-plane traffic generation} through deployed FTD policies and validated the resulting \textbf{connection, prefilter, SSL/TLS, network-discovery, and Security Intelligence events} across \textbf{Inline, Switched, and Routed} interface modes for both \textbf{IPv4 and IPv6} -- 60+ post-migration validation test cases.}
  \resumeItem{Verified \textbf{NAT and Access Control backend} processing post-migration by inspecting rule compilation, installation state, address translation tables, and hit counts on FTD; closed the loop with \textbf{FSIC} (Firepower System Integration Check) and health-alert validation.}
  \resumeItem{Built reusable \textbf{base test classes} shared across the migration suites -- 30+ test cases, 90+ test methods -- and engineered negative-path scenarios stress-testing HA splits, version mismatches, in-flight critical operations, and rollback paths for production firewall appliances.}
\resumeItemListEnd
\vspace{8pt}

\needspace{6\baselineskip}
\textbf{Project 2: HA Failover \& Service-Continuity Validation -- Cisco Secure Firewall Management Center (FMC)}
\nopagebreak
\resumeItemListStart
  \resumeItem{Engineered a \textbf{pyATS} validation framework for FMC \textbf{High Availability} covering \textbf{35+ scenarios and 80+ test methods} -- HA formation, role-switch failover, HA break, sync/mirror control, failure recovery, and negative / degraded-condition paths -- to validate \textbf{network service continuity} after failover events.}
  \resumeItem{Built a reusable \textbf{HA orchestration library} exposing high-level primitives (formation, role-switch, break, sync pause/resume, polling, FTD HA-pair lifecycle) consumed by the broader HA testcase suite.}
  \resumeItem{Layered \textbf{CLI and REST API} probes for device state monitoring, registration verification, and \textbf{managed-device reachability validation} across the \textbf{active/standby network topology} -- confirming policy enforcement and control/management-plane continuity post-failover.}
  \vspace{-4pt}
  \resumeItem{Parallelized execution and refactored fixtures, reducing CI/CD HA-suite runtime from $\sim$6h to $\sim$5h per release cycle.}
\resumeItemListEnd
\vspace{8pt}

\needspace{6\baselineskip}
\textbf{Project 3: Disaster Recovery Automation -- Cloud-Hosted Network Security Management (cdFMC)}
\nopagebreak
\resumeItemListStart
  \resumeItem{Built a \textbf{pre/post state-capture-and-replay DR framework} for cdFMC, automating snapshot lifecycle, recovery orchestration, and post-DR re-validation across distributed cloud-hosted infrastructure.}
  \resumeItem{Authored \textbf{160+ test methods} across four DR scripts -- pre-DR preparation, pre-DR snapshot capture, post-DR validation, and post-DR snapshot re-validation -- covering the full snapshot \& recovery lifecycle.}
  \resumeItem{Re-validated \textbf{14+ configuration and policy domains} -- SRU/LSP versions, registered components, device health, AC / file / identity / SSL-TLS / VPN policies, and deployment history -- to confirm full policy parity and \textbf{device continuity} after recovery.}
  \resumeItem{Operationalized the DR suite inside the cdFMC release \textbf{CI/CD} pipeline, cutting DR validation effort by $\sim$\textbf{90\%} and manual recovery effort by $\sim$\textbf{80\%} per release.}
\resumeItemListEnd
\vspace{8pt}

\needspace{7\baselineskip}
\textbf{Project 4: Backup \& Restore Automation -- Multi-Platform Firewall Devices}
\nopagebreak
\resumeItemListStart
  \resumeItem{Designed a \textbf{pyATS automation framework} spanning \textbf{100+ backup/restore scenarios and 270+ test methods} across on-prem \textbf{FMC}, \textbf{cdFMC}, and \textbf{FTD} devices.}
  \resumeItem{Covered \textbf{four recovery modes} -- standalone, HA, cross-device, and hardware-replacement -- with end-to-end validation of backup integrity, post-restore device reachability, and policy consistency.}
  \resumeItem{Integrated \textbf{multiple storage backends}: Local, \textbf{NFS}, \textbf{SFTP/SCP}, and \textbf{Amazon S3} -- with authentication, transfer validation, and manifest-driven integrity checks.}
  \resumeItem{Built \textbf{REST API test coverage} for on-demand backup operations and the \textbf{Amazon S3 backup/restore} flow across multi-component topologies.}
  \resumeItem{Authored scheduling-automation modules validating cron-driven backup jobs, compressing multi-day manual validation cycles into unattended CI runs.}
\resumeItemListEnd

\vspace{8pt}
\textbf{Recognition: Cisco Internal Hackathon Winner -- KHOJ} \href{https://github.com/Somasekhar9666/resume/blob/main/khoj-hackathon-winner.jpeg}{\small \faTrophy} \hfill \textit{\small Apr 2024}
\nopagebreak
\resumeItemListStart
  \resumeItem{Won \textbf{1st place} as \textbf{sole developer in 72 hours}; built KHOJ, integrating an \textbf{enterprise LLM} with \textbf{MCP-based service orchestration} and a Webex chat front-end.}
\resumeItemListEnd
\resumeSubHeadingListEnd
\vspace{2pt}

%-----------TECHNICAL SKILLS-----------
\section{Technical Skills}
\begin{itemize}[leftmargin=0.15in, label={}]
  \item \textbf{Networking:} TCP/IP, Routing \& Switching fundamentals, NAT, ACLs, Subnetting, VLANs, DNS, DHCP, SSL/TLS, IPv4/IPv6
  \item \textbf{Cisco Security Platforms:} Cisco Secure Firewall (FTD), Firewall Management Center (FMC), Cloud-Hosted FMC (cdFMC), High Availability, Snort3
  \item \textbf{FTD Policy Domains:} Access Control, Prefilter, NAT, SSL/TLS, Security Intelligence, Network Discovery
  \item \textbf{Firewall Topologies:} Inline Sets, Bridge Groups, Routed Interfaces, Static Routes, Security Zones, Port Mode interfaces
  \item \textbf{Network Automation:} pyATS, pytest, Netmiko (SSH-based CLI device automation), REST/JSON API automation, CLI device automation, polling \& retry patterns, fixture design, YAML testbeds
  \item \textbf{Traffic \& Event Validation:} Data-plane traffic generation, event-stream validation (connection / prefilter / SSL-TLS / discovery / SI; IPv4-IPv6), NAT \& AC backend rule inspection, FSIC, health monitoring
  \item \textbf{Languages:} Python (primary), Shell / Bash Scripting, MySQL
  \item \textbf{CI/CD \& Version Control:} Jenkins, pipeline-as-code, parallel execution, release gating, artifact management, structured logging; Git, GitHub (branching, pull requests, code reviews) for automation code management
  \item \textbf{Observability \& Monitoring:} Datadog (dashboard creation \& monitoring), Splunk (log search \& analysis), health monitoring, event-stream validation
  \item \textbf{Cloud \& Infrastructure:} AWS (EC2, S3), GCP, VMware ESXi, Linux, on-prem / bare-metal appliances
  \item \textbf{Protocols \& Storage:} HTTPS, REST over HTTPS, SSH, SFTP / SCP, TCP/IP; Amazon S3, NFS
  \item \textbf{AI Tools:} GitHub Copilot, Cursor, Windsurf, MCP, LLM integration
  \item \textbf{Tools:} Git, Jira, Postman, VS Code
\end{itemize}
\vspace{4pt}

%-----------CERTIFICATIONS---------------
\section{Certifications}
\begin{itemize}[leftmargin=0.2cm, itemsep=1pt]
    \item CCNA -- Cisco Certified Network Associate
    \item Python Programming -- Simplilearn
    \item Command Line in Linux -- Coursera
\end{itemize}
\vspace{6pt}

%-----------EDUCATION---------------
\section{Education}
\resumeSubHeadingListStart
  \resumeSubheading
    {Madanapalle Institute of Technology \& Science}{May 2019 -- May 2023}
    {Bachelor of Technology in Computer Science (GPA: 7.74 / 10.00)}{Madanapalle, India}
\resumeSubHeadingListEnd


\end{document}




